We conducted an extensive simulation study to assess the finite sample performance of our proposed optimal ITR learning method. We examined two data generating mechanisms, an independent failure time simulation setting using exponential distributions and a dependent failure time simulation setting following \cite{fine_proportional_1999}. For both simulation settings, we simulated two causes of failure: priority cause 1 with failure time $T_1$ and cause 2 with failure time $T_2$. We generated $p$ covariates $Z_i$ from independent Uniform(0,1) distributions. %the first column is a row of $1$s for the intercept. 
We assumed two treatments, $A = 0,1$. Let $\beta_1(A)$ and $\beta_2(A)$ denote the $p\times1$ vectors of cause-specific coefficients dependent on treatment assignment $A$. 
\ifnum\prelim=1
\ifnum\jasa=0
The cause 1 and cause 2 failure time hazard coefficients are provided in Table \ref{tab:sim_coeff}.
\else
The PC and competing risk failure time hazard coefficients are provided in the supplementary material.
\fi\fi
The simulations varied by censoring level, treatment assignment mechanisms (independent of covariates vs conditional on covariates), training sample sizes, covariate dimensions, and (for Fine-Gray only) a low vs. high probability mass for the priority cause.  Treatment was allocated using a Bernoulli trial with propensity $\pi(H) = \left[1 + \exp\left\{-Z_i^T \beta_\pi\right\}\right]^{-1}$,
where the propensity coefficient is
$\beta_\pi = (0, -\frac{1}{2} 1_p)^T$ for the observational data setting and $0_{p+1}$ for the randomized trial setting.

We simulated $100$ datasets and compared our method to dtrSurv which refers to the method in \cite{cho_multi-stage_2023}, AIPWE which refers to the method in \cite{he_optimal_2021}, PMCR which refers to the method in \cite{zhou_restricted_2021}, the zero-order model which is our proposed method but assumes no individual variation in treatment effects across patients, and the observed policy. We set the following forest parameters: $n_\text{tree} = 300$, $n_\text{min} = 5$ and $n_\text{minevents} = 2$. We chose $\tol = 0.07$ for the truncated integration criterion $\phi(\cdot)$.%, and $\tol = (0.12, 0.3)$ for the composite probability criterion. For this simulation study, we only include results for the truncated mean criterion. 
%where $2$ is $\sqrt{5}$ rounded to the nearest integer. 
We evaluated methods using a testing dataset with $N_{\text{eval}} = \numeval$ subjects,% without censoring by looking at the distribution of 
assessing $\phi_1(a,Z)$ and $\phi_2(a,Z)$ across simulations. We calculated the OS and PC CIF value functions and averaged the criteria across all subjects per simulation.We note that our proposed method accommodates more than two treatments and causes. 

\ifnum\prelim=1
\ifnum\jasa=0

\newcommand{\tableExample}{
\begin{table}[H]
\caption{CR Example Dataset}
\label{tab:CR_ExampleData}
\centering
\begin{tabular}[t]{ccccccccl}
\hline
ID & $X$ & $\delta$ & $\delta_0$ & $\delta_1$ & $\delta_2$ & $Z$ & $A$ & $\text{Label}_\delta$\\
\hline \hline
1 & 0.041 & 2 & 1 & 0 & 1 & -0.860 & 0 & Cause 2\\
2 & 1.887 & 0 & 0 & 0 & 0 & -0.171 & 1 & Censor\\
3 & 0.065 & 1 & 1 & 1 & 0 & -0.512 & 0 & Cause 1\\
4 & 1.100 & 2 & 1 & 0 & 1 & -1.845 & 0 & Cause 2\\
\hline
\end{tabular}
\end{table}
}

\newcommand{\tableSimsCoeff}{
\begin{table}[H]
\centering
\caption{Cause-specific failure time hazard coefficients and covariate dimensions in Exponential and Fine-Gray simulation scenarios}
\label{tab:sim_coeff}
\renewcommand{\arraystretch}{1.5} % Increase row spacing
\begin{tabular}{cccc}
\toprule
Model &  & Scenario 1 & Scenario 2 \\
 &  & (low dimension)& (high dimension) \\
\midrule
\multirow{5}{*}{Exponential} 
& $p$ & 3 & 6 \\
& $\beta_1^0$ & $(0,-1.5,-1.9,-0.3)^T$ & $(0,-0.5,-1.2,1.5,0.1,0.6)^T$ \\
& $\beta_1^1$ & $(0,-0.6,-1.9,1.5)^T$ & $(0,0.3,-0.4,-0.2,0.9,1.6)^T$ \\
& $\beta_2^0$ & $(0,1,-0.1,0.4)^T$ & $(0,-0.6,1.4,1.5,2,1)^T$ \\
& $\beta_2^1$ & $(0,1.2,0.2,-1)^T$ & $(0,1.5,-2.1,-0.8,0.1,1)^T$ \\
\midrule
\multirow{5}{*}{Fine-Gray} 
& $p$ & 3 & 6 \\
& $\beta_1^0$ & $(0,0.1,0.3)^T$ & $(0,0.9,-0.7,0.6,-0.1,0.5)^T$ \\
& $\beta_1^1$ & $(0,-1.8,-1.5)^T$ & $(0,-1.2,0.1,-0.5,-0.2,-1.2)^T$ \\
& $\beta_2^0$ & $(0,-1.1,-0.3)^T$ & $(0,-0.5,-0.4,-1.1,-0.8,-0.4)^T$ \\
& $\beta_2^1$ & $(0,-0.2,1.2)^T$ & $(0,0.4,-0.3,0.6,0.2,1.1)^T$ \\
\bottomrule
\end{tabular}
\end{table}
}


\newcommand{\tableMeanSD}{%
\begin{table}[ht]
\label{tab:CR_MeanSD}
\centering
\begin{tabular}{rrr}
  \hline
 & means & sds \\ 
  \hline
observed & 138.86 & 2.48 \\ 
  czmk & 128.99 & 26.33 \\ 
  zom & 191.65 & 4.04 \\ 
   \hline
\end{tabular}
\end{table}
}

\newcommand{\tableMeanSDtog}{
\begin{table}[H]
\label{tab:CR_MeanSDtog}
\caption{Mean and SD of critical values across $\numsims$ simulations}
\centering
\begin{tabular}{lccccc}
\hline
    & \multicolumn{2}{c}{Survival} & \multicolumn{2}{c}{Endpoint} & \\
    \hline
    Method  &  \multicolumn{1}{c}{Mean}  &  \multicolumn{1}{c}{SD}  &       \multicolumn{1}{c}{Mean}  &  \multicolumn{1}{c}{SD}  &  \multicolumn{1}{c}{Difference}  \\
    \hline
\input{Chapters/02_Proj1/03_tables/xTable_1}
\hline
\end{tabular}
\end{table}
}

\newcommand{\comparison}{
\begin{table}[h]
\centering
\begin{tabular}{cccccccl}
\hline
Method & A & Endpoint & Criterion & T & C & $\{\pi\}$ &Dataset\\
\hline
\multirow{2}{*}{Proposed method} & \multirow{2}{*}{Finite} & OS & $E(T\wedge\tau),S(t_0)$ & \multirow{2}{*}{NP} & \multirow{2}{*}{I} & \multirow{2}{*}{Flexible} \\
   & & Competing Risk & $E(T_{cif}\wedge\tau),CIF(t_0)$ & & & \\
Yavuz et al (2018) & Finite & CR fr SMART & CIF & NP & I & Flexible (2stage) & Neuroblastoma from RSS\\
Cho et al. (2022) & Finite & Survival & $E(T\wedge\tau),S(t_0)$ & NP & CI & Flexible & ARIC\\
Zhou et al (2021) &  Finite & Main \& Side Effect & &  & I & restricted; linear &HIV from HSSC\\
CPH & Finite & Survival & $E(T); S(t_0)$ & SP & I & -- & -- \\
\hline
\end{tabular}
\caption{Assumptions of existing methods}
\end{table}

OS = Overall Survival \\
NP = Nonparametric \\
SP = Semiparametric \\
C = Censoring Type \\
I = independent \\
CI = conditionally independent\\
HSSC = Health Sciences from South Carolna \\
}


\tableSimsCoeff
\fi
\fi

\subsection{Exponential Setting}
In this simulation setting, we assume $T_1$ is independent from $T_2$, 
where $T_1 \sim \exp(\lambdaone)$, $T_2 \sim \exp(\lambdatwo)$. Since $T_1 \perp T_2$, $T \sim \exp(\lambdaone + \lambdatwo)$.
%Let failure time from cause 1 ($T_1$) follow an exponential distribution with rate $\lambdaone$. Let failure time from cause 2 ($T_2$) follow an exponential distribution with rate $\lambdatwo$. In this setting, we assume $T_1$ and $T_2$ are independent. Thus, by property of independent distributions, the overall survival $T = \min(T_1,T_2)$ follows an exponential distribution with rate $\lambdaone + \lambdatwo$. 
Let $\lambdaoneshort = \lambdaone$ and $\lambdatwoshort = \lambdatwo$. Then for the truncated mean criterion, we have $
\label{eqn:sims_crit_simpleexp}
%int S(t) dt
\phi_1^\mu(t) = \int_0^\tau \survfunc(t)dt = \int_0^\tau \exp(-(\lambdaoneshort+\lambdatwoshort)t)dt$ and $
%int CIF(t) dt
\phi_2^\mu(t) = \int_0^\tau Pr(T \leq t, \epsilon = 1|Z_i) dt = \int_0^\tau \frac{\lambdaoneshort}{\lambdaoneshort + \lambdatwoshort} - e^{-\lambdatwoshort t} + \frac{\lambdatwoshort}{\lambdaoneshort + \lambdatwoshort}\left(e^{-(\lambdaoneshort +\lambdatwoshort) t} \right) dt
%\\ %%S(t_0)
%\phi_1^{\sigma,t_0}(t) &= \survfunc(t_0) = \exp(-(\lambdaoneshort+\lambdatwoshort)t_0)\\
%CIF(t_0)
%\phi_2^{\sigma,t_0}(t) &= Pr(T \leq t_0, \epsilon = 1|Z_i) = \frac{\lambdaoneshort}{\lambdaoneshort + \lambdatwoshort} - e^{-\lambdatwoshort t_0} + \frac{\lambdatwoshort}{\lambdaoneshort + \lambdatwoshort}\left(e^{-(\lambdaoneshort +\lambdatwoshort) t_0} \right)
$. The derivations can be found in the Supplementary Material. Censoring is sampled from exponential distribution such that approximately $15\%$ of the final dataset is censored for low censoring, or approximately $50\%$ for high censoring. 

In Figure \ref{fig:sims_simpleexp_crit1}, we plot the results from our simulation setting, where each dot represents a simulation iteration for each method. The left side shows the mean truncated OS and the right side shows the area under the cause 1 cumulative incidence curve.

Overall, the variance of the truncated mean criterion across simulations for our proposed method is generally smaller than those of other methods. The variance for simulations with $N = 700$ training data is smaller than that for $N = 300$, but there does not appear to be any obvious difference between the two treatment assigning mechanisms. The zero-order model only picks one treatment for the entire sample, thereby having a single regime across samples for the same setting. Thus, in the chance that the ``worse" treatment is selected, it can be harmful to patients. We can see that compared to the zero-order model and observed policy, our proposed method has higher average survival. Compared to the proposed model, both dtrSurv and PMCR have lower average survival and higher average cumulative incidence. APIWE and the observed policy tend to have better CI but has a lower average survival, meaning they are more catered towards the priority cause, at the expense of OS.
\begin{figure}[H]
    \centering
    \includegraphics[width=0.9\textwidth]{Chapters/02_Proj1/03_figures/CR02.SimpleExp_SimulationResults_20240831_crit1.png}
    \caption{Results from exponential setting simulations. (Left) Estimated mean truncated OS time, and (right) the estimated integrated priority cause CIF under each estimated policy. Each dot represents one simulation replicate. The boxplots are the quartile summary of the points and the black dots are the average values of each method. The rows and columns of boxes represent a data generation scenario, study design and sample size.}
	\label{fig:sims_simpleexp_crit1}
\end{figure}

\subsection{Fine-Gray Setting}
\label{sec:finegray}
Following the simulation set-up in \cite{fine_proportional_1999}, we simulate two competing risks for $i = 1,\dots,N$, where the CIF for failures from cause 1 ($\epsilon = 1$) are given by a unit exponential mixture with mass $1-p$ at $\infty$ when covariate values are $0$, and use proportional subdistribution hazards model to obtain CIF for nonzero covariate values:
 $$\Pr(T_i \le t, \epsilon = 1|Z_i) = 1-[1-p\{1-exp(-t)\}]^{\lambdaone}.$$
More details are provided in the supplementary material. We first examine the mean truncated restricted OS time ($\phi^{\mu}_1 = \int^{\tau}_{0} \survfunc(t)dt$), then the mean truncated restricted PC (cause 1) CIF ($\phi^{\mu}_2 = \int^{\tau}_{0} \Pr(T\leq t, \epsilon = 1)dt$). 
The full simulation results are displayed in Figure \ref{fig:sims_finegray_crit1}, where each dot represents a simulation iteration for each method, and the average across simulations is presented. A higher truncated mean survival time (left) and lower area under CIF for priority cause (right) is ideal.
\begin{figure}[H]
    \centering
    \includegraphics[width=1\textwidth]{Chapters/02_Proj1/03_figures/CR02.FineGray_SimulationResults_20240913_crit1.png}%{Chapters/02_Proj1/03_figures/CR02.FineGray_SimulationResults_20240907_crit1.png}
    %Chapters/02_Proj1/03_figures/CR02.FineGray_SimulationResults_20240831_crit1.png}
    \caption{Results from Fine-Gray setting simulations. (Left) Estimated truncated mean overall survival time, and (right) the estimated integrated priority cause cumulative incidence function under each estimated policy. Each dot represents one simulation replicate. The boxplots are the quartile summary of the points and the black dots are the average values of each method. The rows and columns of boxes represent a data generation scenario, study design and sample size.}
	\label{fig:sims_finegray_crit1}
\end{figure}

When $p_{\text{mass}} = 0.2$, our proposed method generally has the same mean times for OS (albeit slightly lower) and slightly lower mean times for CI. We note that while AIPWE and the observed policy sometimes have a lower CIF average (e.g., $5$ covariates and $p_{\text{mass}} = 0.2$), their mean survival time average is also much lower. We also note that in the second row, our method has a mean lower than dtrSurv but higher than AIPWE and PMCR; and the opposite is true for cumulative incidence (lower compared to dtrSurv, but higher compared to AIPWE and PMCR). This is what we would expect, since dtrSurv determines optimal ITRs based on OS, while AIPWE and PMCR determine based on cumulative incidence. When $p_{\text{mass}} = 0.8$, where PC has strong influence on the data, the OS for all methods are similar with low variance across simulations (see Figure \ref{fig:sims_fg_solo} for fixed y-axis scales). However, when comparing the PC cumulative incidence, our model produces the lowest times on average, with the exception of the observed policy (which also has very low OS on average). 

Figure \ref{fig:sims_fg_solo} displays just the setting of low censoring, high sample size, and $5$ covariates where the treatment assignment did not depend on covariates and $p_\text{mass} = 0.8$ from Figure \ref{fig:sims_finegray_crit1}. We can see that our method performs well, with similar survival times to other methods while producing lower mean truncated cumulative incidence for the PC times on average across simulations.

\begin{figure}[H]
    \centering \includegraphics[width=1\textwidth]{Chapters/02_Proj1/03_figures/CR02.solo.FineGray_SimulationResults_20240913_crit1.png}%{Chapters/02_Proj1/03_figures/CR02.solo.FineGray_SimulationResults_20240907_crit1.png}
    \caption{Fine-Gray setting for low censoring, high sample size, high covariates, high PC probability mass for covariate-independent treatment assignment mechanism. (Left) Estimated truncated mean overall survival time, and (right) the estimated integrated priority cause cumulative incidence function under each estimated policy. Each dot represents one simulation replicate. The boxplots are the quartile summary of the points, and the black dots are the average values of each method.}
	\label{fig:sims_fg_solo}
\end{figure}